\documentclass[a4paper]{arrowhead}

\usepackage[yyyymmdd]{datetime}
\usepackage{etoolbox}
\usepackage[utf8]{inputenc}
\usepackage[T1]{fontenc}
\usepackage{multirow}
\usepackage{longtable}
\usepackage{tikz}
\usepackage{adjustbox}
\AtBeginEnvironment{longtable}{\small}
\BeforeBeginEnvironment{tikzpicture}{\begin{adjustbox}{max width=\textwidth}}
\AfterEndEnvironment{tikzpicture}{\end{adjustbox}}
\usetikzlibrary{positioning,shapes.multipart,arrows.meta}

\renewcommand{\dateseparator}{-}
\setlength{\emergencystretch}{3em}
\lstset{language=HTTP,breakatwhitespace=false,columns=fullflexible}

%% Special references
\newcommand{\mref}[2][]{{\textcolor{ArrowheadPurple}{\hyperref[sec:model:#2]{\ifstrempty{#1}{#2}{#1}}}}}
\newcommand{\oref}[1]{{\textcolor{ArrowheadBlue}{\hyperref[sec:operations:#1]{#1}}}}
\newcommand{\pdef}[1]{{\textcolor{ArrowheadGrey}{#1\label{sec:model:primitives:#1}\label{sec:model:primitives:#1s}\label{sec:model:primitives:#1es}}}}
\newcommand{\pref}[1]{{\textcolor{ArrowheadGrey}{\hyperref[sec:model:primitives:#1]{#1}}}}

% as template_sd.tex; the input and output arguments are passed already formatted (\mref or \pref)
\newrobustcmd\osubsection[3]{
  \addtocounter{subsection}{1}
  \addcontentsline{toc}{subsection}{\protect\numberline{\thesubsection}operation \textcolor{ArrowheadBlue}{#1}}
  \subsection*{
    \thesubsection\quad
    operation
    \textcolor{ArrowheadBlue}{#1}
    (\notblank{#2}{#2}{})
    \notblank{#3}{: #3}{}
  }
  \label{sec:operations:#1}
}
\newrobustcmd\msubsection[2]{
  \addtocounter{subsection}{1}
  \addcontentsline{toc}{subsection}{\protect\numberline{\thesubsection}#1 \textcolor{ArrowheadPurple}{#2}}
  \subsection*{\thesubsection\quad#1 \textcolor{ArrowheadPurple}{#2}}
  \label{sec:model:#2} \label{sec:model:#2s} \label{sec:model:#2es}
}
%%

\usepackage{xurl}
\begin{document}
\sloppy

%% Arrowhead Document Properties
\ArrowheadTitle{plantDescriptionMonitor}
\ArrowheadServiceID{plantDescriptionMonitor}
\ArrowheadType{Service Description}
\ArrowheadTypeShort{SD}
\ArrowheadVersion{5.3.0}
\ArrowheadDate{2026-11-09}
\ArrowheadAuthor{Jerker Delsing}
\ArrowheadStatus{REVIEW DRAFT}
\ArrowheadContact{jerker.delsing@ltu.se}
\ArrowheadFooter{\href{www.arrowhead.eu}{www.arrowhead.eu}}
\ArrowheadSetup
%%

%% Front Page
\begin{center}
  \vspace*{1cm}
  \huge{\arrowtitle}

  \vspace*{0.2cm}
  \LARGE{\arrowtype}
  \vspace*{1cm}

  \Large{Service ID: \textit{"\arrowid"}}

  \vspace*{\fill}

  % Front Page Image

  \vspace*{1cm}
  \vspace*{\fill}

  % Front Page Abstract
  \begin{abstract}
    This document describes the \texttt{plant\allowbreak{}Description\allowbreak{}Monitor} service, version 1.0.0, produced by the PlantDescriptionEngine support system of Arrowhead (documentation edition 5.3.0).
  \end{abstract}

  \vspace*{1cm}
\end{center}

\newpage
%%

%% Table of Contents
\tableofcontents
\newpage
%%

\section{Documentation edition and implementation scope}
This is the 5.3.0 documentation edition, dated 2026-11-09. It carries forward the system and service contracts from the retained 5.2.2 documentation and identifies the current primary implementation sources. The former \path|documentation/v5.2.2| folder is retained unchanged.

Document version 5.3.0 is distinct from executable, dependency, registered system/service and wire-contract versions. Values still defined by the source (including 5.2.2 artefact versions, 1.0.0 service versions and \path|SysMonitor-AMQP-IDD-5.2.2| payload identifiers) are preserved. This documentation update does not change runtime behaviour or establish a new implementation release. Historical verification statements describe the tests recorded in the inherited baseline; this revision validates document compilation and references.

\section{Overview}
\label{sec:overview}

This document describes the \texttt{plant\allowbreak{}Description\allowbreak{}Monitor} service, which provides read access to the plant descriptions enriched with the matched registered systems and their monitoring state, and to the alarms raised when the running cloud differs from the description.
It is produced by the PlantDescriptionEngine system \cite{pde-sysd} in service version 1.0.0.

The rest of this document is organized as follows.
In Section \ref{sec:operations}, we describe the abstract message operations provided by the service.
In Section \ref{sec:model}, we end the document by presenting the data types used by the mentioned operations.
The concrete interface is specified in the HTTP/TLS/JSON IDD \cite{pde-plantDescriptionMonitor-http-idd}.

\subsection{Significant Prior Art}

Successor of the Arrowhead 4.6 Plant\allowbreak{}Description\allowbreak{}Engine monitor service; liveness now comes from SysMonitor instead of pinging \texttt{monitorable} services.

\subsection{How This Service Is Meant to Be Used}

Operators and dashboards read descriptions and alarms and acknowledge alarms.

\subsection{Important Delimitations and Dependencies}
\label{sec:delimitations}

Monitoring data are only as recent as SysMonitor's history. With \texttt{service.\allowbreak{}authorization.\allowbreak{}enabled=\allowbreak{}true} every call is checked with Consumer\allowbreak{}Authorization.

\newpage

\section{Service Interface}
\label{sec:operations}

This section describes the interfaces to the \texttt{plant\allowbreak{}Description\allowbreak{}Monitor} service.
In particular, each subsection names an abstract operation, an input type and an output type, in that order.
The input type is named inside parentheses, while the output type is preceded by a colon.
Input and output types are only denoted when accepted or returned, respectively, by the interface in question.

All abstract data types named in this section are defined in Section \ref{sec:model}.

\begin{figure}[ht!]
  \centering
  \begin{tikzpicture}[font=\small]
    \node[draw=ArrowheadBlue, very thick, rectangle split, rectangle split parts=2, align=center,
          rectangle split part fill={ArrowheadBlue!12,white}, text width=9cm, inner sep=6pt] (b)
      {\guillemotleft{}Service\guillemotright{}\\\textbf{plantDescriptionMonitor}
        \nodepart{second} \texttt{query-\allowbreak{}plant-\allowbreak{}descriptions}\\ \texttt{get-\allowbreak{}plant-\allowbreak{}description}\\ \texttt{query-\allowbreak{}alarms}\\ \texttt{get-\allowbreak{}alarm}\\ \texttt{acknowledge-\allowbreak{}alarm}};
  \end{tikzpicture}
  \caption{SysML-style block of the \texttt{plant\allowbreak{}Description\allowbreak{}Monitor} service and its operations.}
  \label{fig:SD}
\end{figure}

The following interface operations are available.

\osubsection{query-plant-descriptions}{}{\mref[Monitor\allowbreak{}Plant\allowbreak{}Description\allowbreak{}List\allowbreak{}Response]{MonitorPlantDescriptionListResponse}}

Returns a page of plant descriptions with monitoring data.

\osubsection{get-plant-description}{\pref{Identifier}}{\mref[Monitor\allowbreak{}Plant\allowbreak{}Description\allowbreak{}Response]{MonitorPlantDescriptionResponse}}

Returns a plant description with monitoring data.

\osubsection{query-alarms}{}{\mref{PdeAlarmListResponse}}

Returns a page of alarms, filtered by system name, severity and acknowledgement.

\osubsection{get-alarm}{\pref{Identifier}}{\mref{PdeAlarmResponse}}

Returns an alarm.

\osubsection{acknowledge-alarm}{\mref{PdeAlarmUpdate}}{\mref{PdeAlarmResponse}}

Acknowledges an alarm.

\newpage

\section{Information Model}
\label{sec:model}

Here, all data objects that can be part of the \texttt{plant\allowbreak{}Description\allowbreak{}Monitor} service are listed in alphabetic order.
Note that each subsection, which describes one type of object, begins with the \textit{struct} keyword, which is used to denote a collection of named fields, each with its own data type.
As a complement to the explicitly defined types in this section, there is also a list of implicit primitive types in Section \ref{sec:model:primitives}, which are used to represent things like hashes and identifiers.

\msubsection{struct}{Connection}

A connection from a consumer port to a producer port.


\noindent\begin{tabularx}{\textwidth}{| >{\raggedright\arraybackslash}p{3.9cm} | >{\raggedright\arraybackslash\small}p{5.3cm} | X |} \hline
\rowcolor{gray!33} Field & Type & Description \\ \hline
consumer & \mref{SystemPort} & The consumer port of the connection. \\ \hline
producer & \mref{SystemPort} & The producer port of the connection. \\ \hline
priority & \pref{Number} & Priority of the resulting orchestration rules (1 highest). \\ \hline
\end{tabularx}



\msubsection{struct}{MonitorPlantDescriptionListResponse}

List of plant descriptions with monitoring data.


\noindent\begin{tabularx}{\textwidth}{| >{\raggedright\arraybackslash}p{3.9cm} | >{\raggedright\arraybackslash\small}p{5.3cm} | X |} \hline
\rowcolor{gray!33} Field & Type & Description \\ \hline
entries & \pref{List}$<$\allowbreak{}\mref[Monitor\allowbreak{}Plant\allowbreak{}Description\allowbreak{}Response]{MonitorPlantDescriptionResponse}$>$ & Matching records of the current page. \\ \hline
count & \pref{Number} & Total number of matching records (all pages). \\ \hline
\end{tabularx}



\msubsection{struct}{MonitorPlantDescriptionResponse}

A plant description with monitoring data.


\noindent\begin{tabularx}{\textwidth}{| >{\raggedright\arraybackslash}p{3.9cm} | >{\raggedright\arraybackslash\small}p{5.3cm} | X |} \hline
\rowcolor{gray!33} Field & Type & Description \\ \hline
id & \pref{Number} & Identifier of the plant description. \\ \hline
plant\allowbreak{}Description & \pref{String} & Name of the plant description. \\ \hline
active & \pref{Boolean} & Whether the entry is in force. \\ \hline
include & \pref{List}$<$\allowbreak{}\pref{Number}$>$ & Identifiers of included plant descriptions. \\ \hline
systems & \pref{List}$<$\allowbreak{}\mref{MonitorSystem}$>$ & Systems to create or update. \\ \hline
connections & \pref{List}$<$\allowbreak{}\mref{Connection}$>$ & Connections between consumer and producer ports. \\ \hline
created\allowbreak{}At & \pref{DateTime} & UTC time of creation. \\ \hline
updated\allowbreak{}At & \pref{DateTime} & UTC time of the last change. \\ \hline
\end{tabularx}



\msubsection{struct}{MonitorSystem}

A plant system with its matched registered system and monitoring data.


\noindent\begin{tabularx}{\textwidth}{| >{\raggedright\arraybackslash}p{3.9cm} | >{\raggedright\arraybackslash\small}p{5.3cm} | X |} \hline
\rowcolor{gray!33} Field & Type & Description \\ \hline
system\allowbreak{}Id & \pref{Identifier} & Identifier of the system within the plant description. \\ \hline
system\allowbreak{}Name & \pref{Name} & Name of the system, PascalCase (e.g. \texttt{TemperatureSensor}). \\ \hline
metadata & \pref{Metadata} & Metadata the registered system must have. \\ \hline
ports & \pref{List}$<$\allowbreak{}\mref{Port}$>$ & Services the system consumes or produces. \\ \hline
registered\allowbreak{}System\allowbreak{}Name & \pref{String} & Name of the registered system matched to this plant system. \\ \hline
monitoring & \pref{Map}$<$\allowbreak{}\pref{String}, \pref{Object}$>$ & Latest monitoring information from SysMonitor. \\ \hline
last\allowbreak{}Observation\allowbreak{}Time & \pref{DateTime} & Time of the latest stored observation of the system. \\ \hline
\end{tabularx}



\msubsection{struct}{PdeAlarmListResponse}

List of alarms.


\noindent\begin{tabularx}{\textwidth}{| >{\raggedright\arraybackslash}p{3.9cm} | >{\raggedright\arraybackslash\small}p{5.3cm} | X |} \hline
\rowcolor{gray!33} Field & Type & Description \\ \hline
entries & \pref{List}$<$\allowbreak{}\mref{PdeAlarmResponse}$>$ & Matching records of the current page. \\ \hline
count & \pref{Number} & Total number of matching records (all pages). \\ \hline
\end{tabularx}



\msubsection{struct}{PdeAlarmResponse}

An alarm.


\noindent\begin{tabularx}{\textwidth}{| >{\raggedright\arraybackslash}p{3.9cm} | >{\raggedright\arraybackslash\small}p{5.3cm} | X |} \hline
\rowcolor{gray!33} Field & Type & Description \\ \hline
id & \pref{Number} & Identifier of the alarm. \\ \hline
system\allowbreak{}Id & \pref{Identifier} & Identifier of the system within the plant description. \\ \hline
system\allowbreak{}Name & \pref{Name} & Name of the system the alarm concerns. \\ \hline
acknowledged & \pref{Boolean} & Whether an operator acknowledged the alarm. \\ \hline
severity & \pref{String} & Severity of the alarm. \\ \hline
cause & \pref{String} & Alarm cause: \texttt{SYSTEM\_\allowbreak{}NOT\_\allowbreak{}REGISTERED}, \texttt{MULTIPLE\_\allowbreak{}MATCHES}, \texttt{SYSTEM\_\allowbreak{}NOT\_\allowbreak{}IN\_\allowbreak{}DESCRIPTION}, \texttt{NOT\_\allowbreak{}MONITORED} or \texttt{NO\_\allowbreak{}RECENT\_\allowbreak{}OBSERVATION}. \\ \hline
description & \pref{String} & Human-readable alarm text. \\ \hline
raised\allowbreak{}At & \pref{DateTime} & Time the alarm was raised. \\ \hline
updated\allowbreak{}At & \pref{DateTime} & UTC time of the last change. \\ \hline
cleared\allowbreak{}At & \pref{DateTime} & Time the alarm was cleared, if cleared. \\ \hline
acknowledged\allowbreak{}At & \pref{DateTime} & Time of the acknowledgement. \\ \hline
\end{tabularx}



\msubsection{struct}{PdeAlarmUpdate}

Acknowledgement of an alarm.


\noindent\begin{tabularx}{\textwidth}{| >{\raggedright\arraybackslash}p{3.9cm} | >{\raggedright\arraybackslash\small}p{5.3cm} | X |} \hline
\rowcolor{gray!33} Field & Type & Description \\ \hline
acknowledged & \pref{Boolean} & Set to true to acknowledge the alarm. \\ \hline
\end{tabularx}



\msubsection{struct}{Port}

A service a plant system consumes or produces.


\noindent\begin{tabularx}{\textwidth}{| >{\raggedright\arraybackslash}p{3.9cm} | >{\raggedright\arraybackslash\small}p{5.3cm} | X |} \hline
\rowcolor{gray!33} Field & Type & Description \\ \hline
port\allowbreak{}Name & \pref{String} & Name of the port within its system. \\ \hline
service\allowbreak{}Definition & \pref{Name} & Service definition name. \\ \hline
service\allowbreak{}Interface & \pref{String} & Required interface template of the port, if any. \\ \hline
metadata & \pref{Metadata} & Metadata the service instance must have. \\ \hline
consumer & \pref{Boolean} & True for a consumed service, false for a produced one. \\ \hline
\end{tabularx}



\msubsection{struct}{SystemPort}

Reference to a port of a plant system.


\noindent\begin{tabularx}{\textwidth}{| >{\raggedright\arraybackslash}p{3.9cm} | >{\raggedright\arraybackslash\small}p{5.3cm} | X |} \hline
\rowcolor{gray!33} Field & Type & Description \\ \hline
system\allowbreak{}Id & \pref{Identifier} & Identifier of the system within the plant description. \\ \hline
port\allowbreak{}Name & \pref{String} & Name of the port within its system. \\ \hline
\end{tabularx}



\subsection{Primitives}
\label{sec:model:primitives}

Types and structures mentioned throughout this document that are assumed to be available to implementations of this service.
The concrete interpretations of each of these types and structures must be provided by any IDD document claiming to implement this service.

\noindent\begin{tabularx}{\textwidth}{| p{4.6cm} | X |} \hline
\rowcolor{gray!33} Type & Description \\ \hline
\pdef{Boolean} & One out of \texttt{true} or \texttt{false}. \\ \hline
\pdef{DateTime} & Pinpoints a specific moment in time (UTC). \\ \hline
\pdef{Identifier} & A string that identifies a record, e.g. a UUID or a composite identifier. \\ \hline
\pdef{List} & An \textit{array} of a known number of items, each having type A. \\ \hline
\pdef{Map} & A set of key-value pairs. \\ \hline
\pdef{Metadata} & A set of key-value pairs with arbitrary, possibly nested values. \\ \hline
\pdef{Name} & A string identifier that is intended to be both human and machine-readable. \\ \hline
\pdef{Number} & An integer number. \\ \hline
\pdef{Object} & Any structured value. \\ \hline
\pdef{String} & A chain of characters. \\ \hline
\end{tabularx}

\newpage

\bibliographystyle{IEEEtran}
\bibliography{bibliography}

\newpage

\section{Revision History}
\subsection{Amendments}

\noindent\begin{tabularx}{\textwidth}{| p{1cm} | p{3cm} | p{2cm} | X | p{4cm} |} \hline
\rowcolor{gray!33} No. & Date & Version & Subject of Amendments & Author \\ \hline

1 & 2026-10-23 & 5.2.2 & Document created from the Arrowhead LaTeX templates for the 5.2.2 Java Spring implementation, including the local extensions of this workspace. & Jerker Delsing \\ \hline
2 & 2026-10-30 & 5.2.2 & Checked against the primary Java source tree (Eclipse-Arrowhead-java): services, operations, paths and data models unchanged. & Claude Opus 5.5 (draft) \\ \hline
3 & 2026-11-09 & 5.3.0 & New documentation edition; primary-source references and Go/Rust implementation scope; 5.2.2 retained. & Documentation maintenance \\ \hline
\end{tabularx}

\subsection{Quality Assurance}

\noindent\begin{tabularx}{\textwidth}{| p{1cm} | p{3cm} | p{2cm} | X |} \hline
\rowcolor{gray!33} No. & Date & Version & Approved by \\ \hline

1 & & \arrowversion  &  \\ \hline

\end{tabularx}

\end{document}
